\documentclass[11pt]{article}

\usepackage[margin=1in]{geometry}
\usepackage{amsmath,amssymb,amsthm}
\usepackage{bm}
\usepackage{mathtools}

% ---- notation macros -------------------------------------------------------
\newcommand{\R}{\mathbb{R}}
\newcommand{\E}{\mathbb{E}}
\newcommand{\KL}{\mathrm{KL}}
\newcommand{\Cat}{\mathrm{Cat}}
\newcommand{\Nor}{\mathcal{N}}
\newcommand{\given}{\,\vert\,}
\newcommand{\sigsq}{\sigma^{2}}
\newcommand{\sigz}[1]{\sigma_{0#1}^{2}}
\newcommand{\const}{\mathrm{const}}
\newcommand{\xim}{x_i^{(m)}}
\newcommand{\half}{\tfrac{1}{2}}
\DeclareMathOperator*{\argmax}{arg\,max}

\theoremstyle{plain}
\newtheorem{proposition}{Proposition}
\theoremstyle{remark}
\newtheorem{remark}{Remark}

\title{\textbf{The per-individual mixture-of-codings single-effect regression}\\[4pt]
\large Mean-field variational EM: model, free energy, and update equations}
\author{William R.\ P.\ Denault}
\date{\today}

\begin{document}
\maketitle

\begin{abstract}
We derive, from first principles, the mean-field variational EM algorithm
implemented in \texttt{ser\_pim\_em} (package \texttt{MixSER}). The model is a
single causal locus at which different individuals may act through different
genetic coding models (additive, dominant, recessive); a latent per-individual
assignment selects the coding, each coding carries one shared random effect, and
a single intercept absorbs the trait baseline. We specify the variational family,
write the evidence lower bound (ELBO / negative free energy) in closed form, and
obtain every update --- the responsibilities, the per-coding effect posterior,
the mixing weights, the intercept, and the empirical-Bayes prior variances --- as
the exact block-coordinate maximiser of that single objective. Consequently the
ELBO is non-decreasing at every step. We close with the two reasons the
implementation uses an intercept rather than mean-centring, a note on
identifiability at a single locus, and an appendix on the analytic single-coding
(Level-1) locus Bayes factor that the routine also returns.
\end{abstract}

% ===========================================================================
\section{Model and notation}\label{sec:model}

Index individuals by $i=1,\dots,n$ and coding models by $m=1,\dots,M$
(typically $M=3$: additive, dominant, recessive). For coding $m$, let
$x^{(m)}=(x_1^{(m)},\dots,x_n^{(m)})^{\!\top}\in\R^{n}$ be the raw genotype
encoding of the single locus under that model (e.g.\ additive $0/1/2$, dominant
$\mathbf 1\{g\ge 1\}$, recessive $\mathbf 1\{g=2\}$). Let $y\in\R^{n}$ be the
response (a partial residual when used inside SuSiE). The residual variance
$\sigsq$ is fixed externally (the outer IBSS loop updates it).

The generative model introduces a latent coding label $z_i\in\{1,\dots,M\}$ for
each individual, one shared effect $\beta_m$ per coding, and a shared intercept
$\alpha$:
\begin{align}
  z_i &\sim \Cat(\bm\pi), &&\bm\pi=(\pi_1,\dots,\pi_M),\ \textstyle\sum_m\pi_m=1,
  \label{eq:gen-z}\\
  \beta_m &\sim \Nor\!\big(0,\ \sigz{m}\big), &&m=1,\dots,M,
  \label{eq:gen-beta}\\
  y_i \given z_i=m,\ \bm\beta &\sim
      \Nor\!\big(\alpha + \xim\,\beta_m,\ \sigsq\big).
  \label{eq:gen-y}
\end{align}
Thus individual $i$ ``reads'' the locus through its latent coding $z_i$, and all
individuals sharing coding $m$ share the \emph{same} effect $\beta_m$. The
unknown parameters are $\bm\theta=(\alpha,\bm\pi,\{\sigz{m}\}_{m})$; the latent
variables are $\bm z=(z_1,\dots,z_n)$ and $\bm\beta=(\beta_1,\dots,\beta_M)$.

The joint density of observed and latent quantities factorises as
\begin{equation}\label{eq:joint}
  p(y,\bm z,\bm\beta\given\bm\theta)
  =\underbrace{\prod_{m=1}^{M}\Nor\!\big(\beta_m;0,\sigz{m}\big)}_{\text{priors on effects}}
   \ \prod_{i=1}^{n}\ \pi_{z_i}\,
     \Nor\!\big(y_i;\,\alpha+x_i^{(z_i)}\beta_{z_i},\,\sigsq\big).
\end{equation}
The quantity of interest, the marginal likelihood
$p(y\given\bm\theta)=\sum_{\bm z}\int p(y,\bm z,\bm\beta\given\bm\theta)\,d\bm\beta$,
is intractable: it couples the $M^{n}$ assignments with the $M$ continuous
effects. We therefore lower-bound it variationally.

% ===========================================================================
\section{Variational family and the free energy}\label{sec:elbo}

\paragraph{Mean-field family.} We approximate the posterior
$p(\bm z,\bm\beta\given y,\bm\theta)$ by a fully factorised distribution
\begin{equation}\label{eq:meanfield}
  q(\bm z,\bm\beta)
  =\Big(\prod_{i=1}^{n} q(z_i)\Big)\Big(\prod_{m=1}^{M} q(\beta_m)\Big),
  \qquad
  q(z_i)=\Cat(\bm r_i),\quad q(\beta_m)=\Nor(\mu_m,\,s_m^{2}).
\end{equation}
Here $r_{im}\coloneqq q(z_i=m)\ge 0$ with $\sum_m r_{im}=1$ are the
\emph{responsibilities}, and $(\mu_m,s_m^{2})$ are the mean and variance of the
Gaussian belief over the coding-$m$ effect. We deliberately keep $q(\beta_m)$
Gaussian; \S\ref{sec:qbeta} shows this family is closed under the update, so the
restriction costs nothing.

\paragraph{ELBO.} For any $q$ the log marginal likelihood admits the
decomposition $\log p(y\given\bm\theta)=\mathcal L(q,\bm\theta)
+\KL\!\big(q\,\Vert\,p(\cdot\given y,\bm\theta)\big)$, where the evidence lower
bound (ELBO; the negative variational free energy) is
\begin{equation}\label{eq:elbo-def}
  \mathcal L(q,\bm\theta)
  =\E_{q}\!\big[\log p(y,\bm z,\bm\beta\given\bm\theta)\big]
   +\mathcal H(q),
  \qquad
  \mathcal H(q)=-\E_q[\log q].
\end{equation}
Because $\KL\ge 0$, $\mathcal L\le\log p(y\given\bm\theta)$, and maximising
$\mathcal L$ jointly over $q$ and $\bm\theta$ is the variational EM problem.

\paragraph{Closed form.} Substituting \eqref{eq:joint} and \eqref{eq:meanfield}
into \eqref{eq:elbo-def} and using $\sum_m r_{im}=1$, the entropy of the
factorised $q$, and $\E_{q(\beta_m)}[\beta_m]=\mu_m$,
$\E_{q(\beta_m)}[\beta_m^2]=\mu_m^2+s_m^2$, gives
\begin{equation}\label{eq:elbo-full}
\boxed{\;
  \mathcal L
  =\sum_{m=1}^{M}\sum_{i=1}^{n} r_{im}\Big(\log\pi_m + \ell_{im}-\log r_{im}\Big)
   \;-\;\sum_{m=1}^{M}\KL\!\big(\Nor(\mu_m,s_m^2)\,\Vert\,\Nor(0,\sigz{m})\big),
\;}
\end{equation}
where the per-pair expected log-likelihood is
\begin{equation}\label{eq:ell}
  \ell_{im}
  \coloneqq \E_{q(\beta_m)}\!\big[\log\Nor(y_i;\alpha+\xim\beta_m,\sigsq)\big]
  =-\half\log(2\pi\sigsq)
   -\frac{1}{2\sigsq}\Big[(y_i-\alpha-\xim\mu_m)^2+(\xim)^2 s_m^2\Big],
\end{equation}
and the Gaussian--Gaussian Kullback--Leibler divergence is the standard
\begin{equation}\label{eq:kl}
  \KL\!\big(\Nor(\mu_m,s_m^2)\,\Vert\,\Nor(0,\sigz{m})\big)
  =\half\!\left[\log\frac{\sigz{m}}{s_m^2}
     +\frac{s_m^2+\mu_m^2}{\sigz{m}}-1\right].
\end{equation}
Equations \eqref{eq:elbo-full}--\eqref{eq:kl} are exactly the ELBO accumulated in
the implementation. The derivation of \eqref{eq:ell} uses only
$\E[(y_i-\alpha-\xim\beta_m)^2]=(y_i-\alpha-\xim\mu_m)^2+(\xim)^2\,\mathrm{Var}_q(\beta_m)$
with $\mathrm{Var}_q(\beta_m)=s_m^2$.

The algorithm maximises the \emph{single} objective \eqref{eq:elbo-full} by
cyclic block-coordinate ascent over the blocks
$\{q(z_i)\}$, $\{q(\beta_m)\}$, $\bm\pi$, $\alpha$, $\{\sigz{m}\}$. We derive each
block's maximiser in turn; since every block is solved exactly, $\mathcal L$
never decreases (\S\ref{sec:monotone}).

% ===========================================================================
\section{E-step: responsibilities $r_{im}$}\label{sec:estep}

Fix $\bm\theta$ and $\{q(\beta_m)\}$ and maximise $\mathcal L$ over a single
$q(z_i)$ subject to $\sum_m r_{im}=1$. Only the first sum in
\eqref{eq:elbo-full} depends on $q(z_i)$:
\begin{equation}
  \mathcal L\big(q(z_i)\big)
  =\sum_{m=1}^{M} r_{im}\big(\log\pi_m+\ell_{im}-\log r_{im}\big)+\const .
\end{equation}
Introducing a Lagrange multiplier $\lambda$ for the constraint and setting
$\partial/\partial r_{im}=0$,
\begin{equation}
  \log\pi_m+\ell_{im}-\log r_{im}-1-\lambda=0
  \quad\Longrightarrow\quad
  r_{im}\propto \pi_m\,e^{\ell_{im}} .
\end{equation}
This is the generic mean-field fixed point
$q^\star(z_i)\propto\exp\E_{q(\bm\beta)}[\log p(y,\bm z,\bm\beta)]$. Dropping the
constant $-\half\log(2\pi\sigsq)$ from \eqref{eq:ell} (it cancels in the
normalisation) yields the implemented form
\begin{equation}\label{eq:rstep}
\boxed{\;
  r_{im}
  =\frac{\pi_m\,\exp(\eta_{im})}{\sum_{k=1}^{M}\pi_k\,\exp(\eta_{ik})},
  \qquad
  \eta_{im}=-\frac{1}{2\sigsq}\Big[(y_i-\alpha-\xim\mu_m)^2+(\xim)^2 s_m^2\Big].
\;}
\end{equation}
The effect is integrated against its \emph{current} belief $q(\beta_m)$: the
score $\eta_{im}$ compares individual $i$'s observation to the fitted value
$\alpha+\xim\mu_m$ (never to $0$), penalised by the predictive uncertainty
$(\xim)^2 s_m^2$. Individuals are thus assigned to codings by fit quality, with a
built-in Occam penalty for uncertain effects.

\begin{remark}[Reference genotype is uninformative]\label{rmk:g0}
For an individual with $g_i=0$ every coding gives $\xim=0$, so
$\eta_{im}=-(y_i-\alpha)^2/(2\sigsq)$ is the \emph{same} for all $m$ and
$r_{im}=\pi_m$. Such individuals carry no coding information and correctly revert
to the prior. This is exactly why the baseline must enter as the intercept
$\alpha$ and not by mean-centring $y$ (\S\ref{sec:intercept}).
\end{remark}

% ===========================================================================
\section{M-step I: per-coding effect posterior $q(\beta_m)$}\label{sec:qbeta}

Fix $\{r_{im}\}$ and $\bm\theta$ and optimise over $q(\beta_m)$. The generic
mean-field optimum is
$\log q^\star(\beta_m)=\E_{q(\bm z)}\big[\log p(y,\bm z,\bm\beta)\big]+\const$,
where the expectation is over everything except $\beta_m$. Collecting the
$\beta_m$-dependent terms of \eqref{eq:joint},
\begin{equation}\label{eq:logqbeta}
  \log q^\star(\beta_m)
  =-\frac{1}{2\sigsq}\sum_{i=1}^{n} r_{im}\big(y_i-\alpha-\xim\beta_m\big)^2
   -\frac{\beta_m^2}{2\,\sigz{m}}+\const .
\end{equation}
This is quadratic in $\beta_m$, so $q^\star(\beta_m)$ is Gaussian (the family is
closed). Matching the coefficients of $\beta_m^2$ and $\beta_m$ gives the
precision and mean:
\begin{align}
  \frac{1}{s_m^2}
  &=\frac{1}{\sigz{m}}+\frac{1}{\sigsq}\sum_{i=1}^{n} r_{im}\,(\xim)^2,
  \label{eq:s2step}\\[4pt]
  \mu_m
  &=\frac{s_m^2}{\sigsq}\sum_{i=1}^{n} r_{im}\,\xim\,(y_i-\alpha).
  \label{eq:mustep}
\end{align}
Writing $S_m\coloneqq\sum_i r_{im}(\xim)^2$, these are the boxed updates
\begin{equation}
\boxed{\;
  s_m^2=\Big(\tfrac{1}{\sigz{m}}+\tfrac{S_m}{\sigsq}\Big)^{\!-1},
  \qquad
  \mu_m=\frac{s_m^2}{\sigsq}\sum_{i=1}^{n} r_{im}\,\xim\,(y_i-\alpha).
\;}
\end{equation}
This is precisely a responsibility-weighted Bayesian single-effect regression of
$y-\alpha$ on coding $m$: each individual contributes with weight $r_{im}$, and
the prior $\Nor(0,\sigz{m})$ shrinks the estimate.

% ===========================================================================
\section{M-step II: hyperparameters}\label{sec:hyper}

\subsection{Mixing weights $\bm\pi$}
Only $\sum_{m}\big(\sum_i r_{im}\big)\log\pi_m$ in \eqref{eq:elbo-full} depends on
$\bm\pi$. Maximising under $\sum_m\pi_m=1$ with a Lagrange multiplier gives the
ordinary EM weight update
\begin{equation}\label{eq:pistep}
\boxed{\;
  \pi_m=\frac{1}{n}\sum_{i=1}^{n} r_{im}.
\;}
\end{equation}
No \texttt{mixsqp} / empirical-Bayes step is applied to the per-individual
responsibilities; doing so collapses the weights onto the single sparsest coding
(see \S\ref{sec:intercept}).

\subsection{Intercept $\alpha$}
Differentiating \eqref{eq:elbo-full} through \eqref{eq:ell},
\begin{equation}
  \frac{\partial\mathcal L}{\partial\alpha}
  =\frac{1}{\sigsq}\sum_{i=1}^{n}\sum_{m=1}^{M} r_{im}\big(y_i-\alpha-\xim\mu_m\big)=0 .
\end{equation}
Using $\sum_m r_{im}=1$ and solving,
\begin{equation}\label{eq:alphastep}
\boxed{\;
  \alpha=\frac{1}{n}\sum_{i=1}^{n}\Big(y_i-\sum_{m=1}^{M} r_{im}\,\xim\,\mu_m\Big).
\;}
\end{equation}
The intercept is the mean residual after removing the responsibility-averaged
mixed effect.

\subsection{Empirical-Bayes prior variances $\sigz{m}$}
The only $\sigz{m}$-dependent part of \eqref{eq:elbo-full} is
$-\KL(\Nor(\mu_m,s_m^2)\Vert\Nor(0,\sigz{m}))$. Holding $q(\beta_m)$ fixed and
writing $v\coloneqq\sigz{m}$,
\begin{equation}
  \frac{\partial}{\partial v}\Big[-\half\log v-\frac{s_m^2+\mu_m^2}{2v}\Big]
  =-\frac{1}{2v}+\frac{s_m^2+\mu_m^2}{2v^2}=0
  \;\Longrightarrow\;
  v=s_m^2+\mu_m^2 .
\end{equation}
Hence the ELBO-maximising empirical-Bayes update is the second moment of the
effect belief,
\begin{equation}\label{eq:ebstep}
\boxed{\;
  \sigz{m}=\mu_m^2+s_m^2=\E_{q(\beta_m)}[\beta_m^2]
\;}
\end{equation}
(applied when \texttt{estimate\_prior\_variance=TRUE}; a floor $10^{-8}$ keeps it
positive). Equations \eqref{eq:s2step} and \eqref{eq:ebstep} are coupled across
iterations but each is the exact maximiser of its own block, which is all
coordinate ascent requires.

% ===========================================================================
\section{Monotonicity, the algorithm, and convergence}\label{sec:monotone}

\begin{proposition}[Monotone ELBO]
Each update \eqref{eq:rstep}, \eqref{eq:s2step}--\eqref{eq:mustep},
\eqref{eq:pistep}, \eqref{eq:alphastep}, \eqref{eq:ebstep} is the exact global
maximiser of the single objective $\mathcal L$ in \eqref{eq:elbo-full} over its
block, holding the others fixed. Therefore the sequence of ELBO values produced
by cyclic application is non-decreasing,
$\mathcal L^{(t+1)}\ge\mathcal L^{(t)}$, and, being bounded above by
$\log p(y\given\bm\theta)$, it converges.
\end{proposition}
\begin{proof}
Block-coordinate ascent on a function with exact block maximisers is monotone:
replacing one block by its argmax cannot decrease the objective, and the other
blocks are unchanged. The $q(z_i)$ and $q(\beta_m)$ steps are the mean-field
optima (\S\ref{sec:estep},\S\ref{sec:qbeta}); the $\bm\pi$, $\alpha$, $\sigz{m}$
steps are the stationary points of concave one-block subproblems
(\S\ref{sec:hyper}). The bound $\mathcal L\le\log p(y\given\bm\theta)$ follows
from $\KL\ge 0$. A monotone sequence bounded above converges.
\end{proof}

\paragraph{Algorithm (one sweep).} Given $\sigsq$, repeat to convergence of
$\mathcal L$:
\begin{enumerate}\itemsep2pt
  \item \textbf{E-step:} update all $r_{im}$ by \eqref{eq:rstep};
  \item \textbf{weights:} $\pi_m\leftarrow\frac1n\sum_i r_{im}$
        \eqref{eq:pistep};
  \item \textbf{intercept:} $\alpha$ by \eqref{eq:alphastep};
  \item \textbf{effect:} $s_m^2,\mu_m$ by \eqref{eq:s2step}--\eqref{eq:mustep};
  \item \textbf{EB prior:} $\sigz{m}\leftarrow\mu_m^2+s_m^2$ \eqref{eq:ebstep}.
\end{enumerate}
Initialise $\bm\pi$ uniform, $\alpha=\bar y$, and each $\mu_m$ at the single-coding
OLS slope of $y-\alpha$ on $x^{(m)}$ with $s_m^2=\sigsq/\|x^{(m)}\|^2$. Because
the per-individual mixture is only weakly identified at one locus
(\S\ref{sec:intercept}), the routine runs \texttt{n\_start} restarts --- the
first uniform, the rest $\mathrm{Dirichlet}(\mathbf 1)$ draws of $\bm\pi$ --- and
returns the fit of highest ELBO.

\paragraph{Variational marginal likelihood.} A convenient by-product evaluated
during the E-step is
\begin{equation}\label{eq:loglik}
  \widetilde{\ell}(y)
  =\sum_{i=1}^{n}\log\sum_{m=1}^{M}\pi_m\,e^{\ell_{im}},
\end{equation}
the log of the (effect-integrated) per-individual mixture likelihood at the
current beliefs; it is reported as \texttt{loglik}.

% ===========================================================================
\section{Why an intercept, not mean-centring; identifiability}
\label{sec:intercept}

\paragraph{Intercept vs.\ centring.} By Remark~\ref{rmk:g0}, raw $0/1/2$ codings
make the reference genotype ($g_i=0$) predict $\alpha$ under every coding, so such
individuals are uninformative ($r_{im}=\pi_m$). If instead one mean-centres $y$
and drops $\alpha$, the $g=0$ group sits at $y\approx-\bar y$ while the model
predicts $0$ there; every component must then bend to absorb the baseline, and
the optimiser collapses onto whichever coding can do so most cheaply --- the
spurious ``recessive collapse''. The shared intercept \eqref{eq:alphastep}
removes the baseline without this distortion and is invariant to the trait mean.

\paragraph{Weights by EM, not \texttt{mixsqp}.} Treating the columns
$\{r_{\cdot m}\}$ as data for an empirical-Bayes mixing solver
(\texttt{mixsqp}) maximises a different objective that rewards sparsity in
$\bm\pi$ and again collapses to one coding. The plain EM update
\eqref{eq:pistep} is the correct M-step for \eqref{eq:elbo-full}.

\paragraph{Identifiability at one locus.} The dominant coding
$\mathbf 1\{g\ge1\}$ equals the additive coding except at $g=2$, and the
recessive coding $\mathbf 1\{g=2\}$ is informative only at $g=2$; the
contrasts that separate codings live entirely in the $g\ge 1$ stratum, and
$g=0$ individuals are uninformative. A single locus therefore identifies the
per-individual mixture only weakly, and a residual dominant component can
persist. The restarts guard the optimisation, but \emph{calibrated} weights
$\bm\pi$ require pooling the responsibilities across many loci; at one locus the
quantity to trust for a decision is the analytic Level-1 Bayes factor of
Appendix~\ref{app:bf}, not $\bm\pi$ itself.

% ===========================================================================
\appendix
\section{Analytic single-coding (Level-1) locus Bayes factor}\label{app:bf}

Independently of the per-individual mixture, \texttt{ser\_pim\_em} also returns
the closed-form Bayes factor of ``some coding has a non-zero effect'' versus
``all effects are zero''. For each coding $m$ let $\tilde x^{(m)}$ be the
centred-and-scaled $x^{(m)}$ and $\tilde y=y-\bar y$. The exact single-effect
Bayes factor of Wang et al.\ (2020) for coding $m$ at prior variance $V$ is
\begin{equation}\label{eq:bf1}
  \log\mathrm{BF}_m(V)
  =\half\log\frac{\hat s_m^2}{V+\hat s_m^2}
   +\frac{\hat\beta_m^2}{2}\Big(\frac{1}{\hat s_m^2}-\frac{1}{V+\hat s_m^2}\Big),
  \quad
  \hat\beta_m=\frac{\tilde x^{(m)\top}\tilde y}{\|\tilde x^{(m)}\|^2},
  \quad
  \hat s_m^2=\frac{\sigsq}{\|\tilde x^{(m)}\|^2}.
\end{equation}
Each $V_m$ is set by maximum marginal likelihood,
$V_m=\argmax_{V\ge 0}\log\mathrm{BF}_m(V)$ (with the off-switch $V_m=0$, i.e.\
$\log\mathrm{BF}_m=0$, whenever the optimum is negative). With coding weights
$w_m$ (uniform $1/M$ by default) the locus Bayes factor of any-coding-effect vs.\
null is the weighted average
\begin{equation}\label{eq:bfsnp}
  \log\mathrm{BF}_{\mathrm{locus}}
  =\log\sum_{m=1}^{M} w_m\,\mathrm{BF}_m(V_m)
  =\operatorname*{logsumexp}_{m}\!\big(\log w_m+\log\mathrm{BF}_m(V_m)\big),
\end{equation}
returned as \texttt{lbf\_snp}/\texttt{bf\_snp}. This is the object the SuSiE /
SuSiE~2.0 IBSS loop consumes as the single-effect log-Bayes-factor for this
locus; the per-individual mixture's exact Bayes factor would require the
intractable $M^{n}$ sum, so for decisions the analytic \eqref{eq:bfsnp} is used.
As a conservative cross-check the routine also reports the variational lower
bound $\log\mathrm{BF}_{\mathrm{pim}}=\mathcal L-\ell_0$, where
$\ell_0=\sum_i\log\Nor(y_i;\bar y,\sigsq)$ is the intercept-only null
log-likelihood.

\paragraph{Reference.}
G.\ Wang, A.\ Sarkar, P.\ Carbonetto \& M.\ Stephens (2020).
\textit{A simple new approach to variable selection in regression, with
application to genetic fine mapping.} JRSS-B 82(5), 1273--1300.

\end{document}
